\chapter{Abel's Theorem and the Jacobian}
\label{ch:abel-jacobian}

Let \(T=\C/\Lambda\).  If
\[
  D=\sum n_i[p_i]\in\Div^0(T),
\]
choose lifts \(z_i\in\C\) of \(p_i\) and define
\[
  \smap(D)=\sum n_i z_i \pmod{\Lambda}.
\]
Because \(\sum n_i=0\), changing lifts changes the sum by an element of
\(\Lambda\).  Thus \(\smap(D)\in T\) is well-defined.

\begin{theorem}[Abel's theorem on a torus]
\label{thm:abel-torus-new}
A degree-zero divisor \(D\) on \(T\) is principal if and only if
\(\smap(D)=0\) in \(T\).
\end{theorem}

\begin{proof}[Proof sketch]
If \(D=\divisor(f)\), the elliptic Liouville laws of
\cref{thm:elliptic-liouville} show that the Abel sum vanishes.  Conversely, if
the degree and Abel sum vanish, one constructs an elliptic function with exactly
the prescribed zeros and poles.  A standard construction uses the Weierstrass
\(\sigma\)-function; Riemann--Roch gives another route.
\end{proof}

\begin{corollary}
\label{cor:pic-torus}
The Abel--Jacobi map induces an isomorphism
\[
  \Pic^0(T)\cong T=\C/\Lambda.
\]
\end{corollary}

Transporting this statement to the cubic gives the geometric identification
\[
  \Pic^0(E)\cong E(\C).
\]
The map is \(P\mapsto[P]-[O]\), and Abel's theorem tells us exactly when two such
degree-zero divisor classes are equal.

\section*{Exercises}
\addcontentsline{toc}{section}{Exercises}

\exerciseintent{Compute the Abel sum of the simplest degree-zero divisor.}
\begin{exercise}
Choose lifts \(p,q\in\C\) of \(P,Q\in T\).  Compute
\(\smap([P]-[Q])\) from the definition and show that it is the class of
\(p-q\) modulo \(\Lambda\), hence \(P-Q\) in the group \(T\).
\end{exercise}

\exerciseintent{Use Abel's theorem to recognize when a two-point divisor is principal.}
\begin{exercise}
Choose lifts \(p,q,0\in\C\) of \(P,Q,O\).  Compute the Abel sum of
\([P]+[Q]-2[O]\).  Use Abel's theorem to turn the congruence
\(p+q\equiv0\pmod{\Lambda}\) into a principal-divisor criterion.
\end{exercise}

\exerciseintent{Use Abel sums to prove injectivity of the point-to-divisor-class map.}
\begin{exercise}
Assume \([P]-[O]=[Q]-[O]\) in \(\Pic^0(T)\).  Subtract the two divisor classes,
compute the Abel sum of \([P]-[Q]\), and conclude that \(P=Q\).
\end{exercise}

\exerciseintent{Check additivity of the Abel-sum map term by term.}
\begin{exercise}
Write
\[
  D=\sum_i n_i[p_i],\qquad D'=\sum_j m_j[q_j],
\]
with chosen lifts of all points.  Compute both sides of
\[
  \smap(D+D')=\smap(D)+\smap(D')
\]
term by term in \(\C/\Lambda\).
\end{exercise}

\exerciseintent{Verify that principal divisors are killed so the Abel-sum map descends to \(\Pic^0\).}
\begin{exercise}
Let \(D\in\Div^0(T)\) and let \(f\) be a nonzero meromorphic function on \(T\).
Compute
\[
  \smap(D+\divisor(f))-\smap(D).
\]
Use Abel's theorem to show that this difference is \(0\), and then verify that
the induced map \(\Pic^0(T)\to T\) is a group homomorphism.
\end{exercise}
